\chapter{Géométrie complexe et racines de l’unité}
\section*{Exercice 28.1 -- Distances}
Le premier est le cercle de centre d'affixe $1$ et rayon $2$. Pour le second, écrire $z=x+i y$ donne $(x-1)^2+y^2=(x+1)^2+y^2$, donc $x=0$ : l'axe imaginaire, médiatrice du segment d'extrémités $-1,1$.

\section*{Exercice 28.2 -- Rotation}
Elle s'écrit $z'-1=i(z-1)$. C'est la rotation de centre d'affixe $1$ et d'angle $\pi/2$ modulo $2\pi$. Elle conserve les distances car $|i|=1$.

\section*{Exercice 28.3 -- Angle droit}
$(a-c)/(b-c)=(-1-i)/(1-i)=-i$, imaginaire pur non nul. Son argument est $-\pi/2$, donc les vecteurs $\overrightarrow{CB}$ et $\overrightarrow{CA}$ sont perpendiculaires. On peut aussi vérifier $AC^2+BC^2=2+2=4=AB^2$.

\section*{Exercice 28.4 -- Racines quatrièmes}
Les solutions sont $e^{2ki\pi/4}$ pour $k=0,1,2,3$, soit $1,i,-1,-i$. Leur somme vaut zéro. Elles sont les sommets d'un carré centré à l'origine.

\section*{Exercice 28.5 -- Racines d’un nombre non nul}
$-8i=8e^{-i\pi/2}$. Chaque solution a module $2$ et un argument $-\pi/6+2k\pi/3$, $k=0,1,2$. Les solutions sont donc $\sqrt3-i$, $2i$, $-\sqrt3-i$. Leur cube vaut bien le nombre demandé.

\section*{Exercice 28.6 -- Somme des racines de l’unité}
Avec $q=e^{2i\pi/n}\ne1$, la somme est $(1-q^n)/(1-q)=0$. Le barycentre à poids égaux des sommets du polygone régulier a donc affixe zéro, son centre. Pour $n=1$, la somme serait $1$, d'où l'hypothèse.
